\documentclass[a4paper,12pt]{article}

\usepackage[T2A]{fontenc}
\usepackage[utf8]{inputenc}
\usepackage[russian]{babel}
\usepackage{amsmath,amssymb,mathtools}
\renewcommand{\familydefault}{\sfdefault}
\usepackage{icomma}
\usepackage[a4paper,margin=24mm,headheight=15pt,footskip=15mm]{geometry}
\usepackage{microtype}
\usepackage{tabularx,booktabs,array}
\usepackage{tikz}
\usepackage{fancyhdr}
\usepackage{hyperref}
\usepackage{xcolor}
\usepackage{parskip}
\usepackage{setspace}
\usepackage{needspace}

\definecolor{accent}{HTML}{008080}
\definecolor{darkaccent}{HTML}{004D4D}
\definecolor{textgray}{HTML}{333333}
\definecolor{muted}{HTML}{647576}
\definecolor{reviewred}{HTML}{A33333}
\definecolor{reviewblue}{HTML}{2A607F}

\hypersetup{colorlinks=true,linkcolor=darkaccent,urlcolor=darkaccent}
\setstretch{1.22}
\setlength{\parindent}{0pt}
\setlength{\parskip}{6pt}
\setcounter{secnumdepth}{0}
\renewcommand{\arraystretch}{1.3}
\renewcommand{\headrulewidth}{0.3pt}
\setlength{\headsep}{9mm}
\pagestyle{fancy}
\fancyhf{}
\fancyhead[L]{\footnotesize\textcolor{muted}{Отношения, пропорции и формулы}}
\fancyhead[R]{\footnotesize\textcolor{muted}{Надя Клименко}}
\fancyfoot[C]{\footnotesize\thepage}
\fancypagestyle{plain}{\fancyhf{}\fancyfoot[C]{\footnotesize\thepage}\renewcommand{\headrulewidth}{0pt}}

\newcommand{\accentline}{\par\noindent\textcolor{accent}{\rule{\linewidth}{0.45pt}}\par}
\newcommand{\reviewheading}[1]{\Needspace{8\baselineskip}\par\vspace{0.75em}{\large\bfseries\textcolor{darkaccent}{#1}}\par\accentline}
\newcommand{\field}[1]{\textbf{#1}}

\begin{document}
\pagenumbering{roman}

\begin{titlepage}
  \thispagestyle{empty}
  \begin{center}
    \vspace*{36mm}
    {\small\textcolor{muted}{МАТЕМАТИКА}\par}
    \vspace{24mm}
    {\Huge\bfseries\textcolor{darkaccent}{Ревью домашней работы}\par}
    \vspace{17mm}
    {\Large Отношения, пропорции и преобразование формул\par}
    \vspace{32mm}
    {\large Надя Клименко\par}
    \vspace{18mm}
    {\normalsize 23 ноября 2026 г.\par}
    \vfill
    {\small Автор: Лёвин Артём Александрович\par}
    \vspace{3mm}
    {\small\textcolor{muted}{Лёвин Артём Александрович эксклюзивно для Нади Клименко}\par}
    \vspace{12mm}
    \begin{tikzpicture}[x=1cm,y=1cm]
      \draw[accent,line width=0.7pt]
        (-4.6,0) .. controls (-2.5,0.65) and (-1.0,-0.4) .. (0,0)
                 .. controls (1.2,0.35) and (3.0,-0.55) .. (4.6,0);
    \end{tikzpicture}
  \end{center}
\end{titlepage}
\pagenumbering{arabic}

\clearpage
\section*{\textcolor{darkaccent}{Сводная таблица}}
\accentline

В таблицу включены только задания, требующие исправления. Нажмите на номер,
чтобы перейти к подробному разбору.

\small
\begin{tabularx}{\linewidth}{@{}p{12mm}p{28mm}X p{28mm}p{31mm}@{}}
\toprule
\textbf{№} & \textbf{Тема} & \textbf{Суть неточности} &
\textbf{Ваш ответ} & \textbf{Правильный ответ}\\
\midrule
\hyperref[task:10]{10} & Выражение величины
из формулы &
Не указано, что \(r\ne0\); условие \(F\ne0\) добавлено без основания. &
\(F\ne0,\ k\ne0,\ m_1\ne0\) &
\(r\ne0,\ k\ne0,\ m_1\ne0\)\\
\bottomrule
\end{tabularx}
\normalsize

\vspace{9mm}
\textcolor{reviewblue}{\textbf{Главный вывод.}}
При преобразовании буквенной формулы нужно помнить сразу о двух источниках
ограничений: о знаменателях исходного выражения и о величинах, на которые
приходится делить в ходе решения.

\clearpage
\thispagestyle{plain}
\vspace*{\fill}
\begin{center}
  {\Huge\bfseries\textcolor{darkaccent}{Разбор ошибок}\par}
  \vspace{9mm}
  \textcolor{accent}{\rule{65mm}{0.7pt}}
\end{center}
\vspace*{\fill}
\clearpage

\phantomsection\label{task:10}
\section*{\textcolor{darkaccent}{Задание №\,10. Выражение величины из формулы}}
\accentline

\field{Условие.}
Из формулы
\[
  F=\frac{k m_1 m_2}{r^2}
\]
выразите \(m_2\). Укажите, какие величины при этом должны быть ненулевыми.

\reviewheading{Ваше решение}

Последовательность преобразований записана верно:
\[
\begin{aligned}
F&=\frac{k m_1 m_2}{r^2},\\
k m_1 m_2&=F r^2,\\
m_2&=\frac{F r^2}{k m_1}.
\end{aligned}
\]
После этого указаны условия \(F\ne0\), \(k\ne0\), \(m_1\ne0\).

\field{Ваш ответ:}
\[
m_2=\frac{F r^2}{k m_1},\qquad F\ne0,\quad k\ne0,\quad m_1\ne0.
\]

\reviewheading{Ошибка и её точное место}

\textcolor{reviewred}{\textbf{Ошибка:}} неверный перечень ограничений
при правильной формуле для \(m_2\).

\field{Последний правильный шаг:}
\[
  m_2=\frac{F r^2}{k m_1}.
\]

\field{Первый неточный шаг:} перечень \(F\ne0,\ k\ne0,\ m_1\ne0\).
В нём пропущено обязательное условие \(r\ne0\), зато добавлено
необязательное условие \(F\ne0\).

\reviewheading{Почему это неверно}

\textbf{1. Ограничение на \(r\) возникает уже в исходной формуле.}

В знаменателе стоит \(r^2\). Дробь имеет смысл только при ненулевом
знаменателе, поэтому
\[
  r^2\ne0 \quad\Longleftrightarrow\quad \boxed{r\ne0}.
\]
Даже если в итоговой записи \(r^2\) переместилось в числитель,
исходное ограничение сохраняется.
При \(r=0\) исходная формула вообще не определена.

\textbf{2. Ограничения на \(k\) и \(m_1\) связаны с делением.}

Чтобы получить \(m_2\), необходимо разделить обе части равенства
\(k m_1 m_2=F r^2\) на произведение \(k m_1\):
\[
m_2=\frac{F r^2}{k m_1}.
\]
Делить на нуль нельзя, следовательно,
\[
  k m_1\ne0
  \quad\Longleftrightarrow\quad
  \boxed{k\ne0,\qquad m_1\ne0}.
\]
Если \(k=0\) или \(m_1=0\), исходная формула при \(r\ne0\)
даёт \(F=0\), независимо от \(m_2\). В таком случае она не позволяет
однозначно найти \(m_2\).

\textbf{3. Ограничение \(F\ne0\) лишнее.}

В ходе решения мы нигде не делим на \(F\). Поэтому \(F\) вправе
равняться нулю. Например, при
\[
  k=1,\qquad m_1=1,\qquad r=1,\qquad m_2=0
\]
исходная формула даёт
\[
  F=\frac{1\cdot1\cdot0}{1^2}=0.
\]
Все действия при этом допустимы. Требование \(F\ne0\)
необоснованно исключило бы этот правильный случай.

\clearpage
\reviewheading{Ключевая идея}

\textcolor{reviewblue}{\textbf{Ключевая идея:}}
при выражении неизвестной величины проверяйте два вопроса:
какие знаменатели обязаны оставаться ненулевыми в исходной формуле
и на какие множители Вы делите при преобразованиях.

\reviewheading{Исправление от последнего правильного шага}

Исходное выражение для \(m_2\) сохраняем без изменений:
\[
  m_2=\frac{F r^2}{k m_1}.
\]
Исправляем только условия:
\[
  \boxed{r\ne0,\qquad k\ne0,\qquad m_1\ne0}.
\]
Условие \(F\ne0\) убираем.

\reviewheading{Полное правильное решение}

\textbf{Шаг 1.} Исходное равенство:
\[
  F=\frac{k m_1 m_2}{r^2},\qquad r\ne0.
\]

\textbf{Шаг 2.} Умножим обе части на \(r^2\):
\[
  F r^2=k m_1 m_2.
\]

\textbf{Шаг 3.} Разделим обе части на \(k m_1\).
Это допустимо при \(k\ne0\) и \(m_1\ne0\):
\[
  \frac{F r^2}{k m_1}=m_2.
\]

\textbf{Шаг 4.} Запишем ответ с полным перечнем ограничений:
\[
\boxed{
  m_2=\frac{F r^2}{k m_1},
  \qquad r\ne0,\ k\ne0,\ m_1\ne0.
}
\]

\reviewheading{Проверка}

Подставим найденное выражение обратно в правую часть
исходной формулы:
\[
\frac{k m_1 m_2}{r^2}
=
\frac{k m_1}{r^2}\cdot\frac{F r^2}{k m_1}
=F.
\]
Сокращение допустимо при трёх указанных ограничениях.
Следовательно, преобразование выполнено правильно.

\textcolor{darkaccent}{\textbf{Итог.}}
Формула для \(m_2\) у Вас получилась правильная.
Для законченного ответа необходимо указать
\(r\ne0\), \(k\ne0\), \(m_1\ne0\). Значение \(F=0\) допустимо.

\reviewheading{Задача на закрепление}

Из формулы
\[
  T=\frac{6uv}{w^2}
\]
выразите \(v\). Укажите, какие величины при этом должны быть ненулевыми.

\vfill
\begin{center}
\footnotesize\textcolor{muted}{Лёвин Артём Александрович эксклюзивно для Нади Клименко}
\end{center}

\end{document}
